\section*{Acknowledgements}
We thank Jean Zinn--Justin for helpful discussions on
the renormalization of the Langevin equation and for encouraging us
to proceed directly with the flow equation rather than considering
the latter to be a limit of the Langevin equation.

\appendix

\section{Notational conventions}\label{app:A}

The Lie algebra $\sun$ of $\SUn$
may be identified with the linear space of all anti-hermitian
traceless $N\times N$ matrices.
With respect to a
basis $T^a$, $a=1,\ldots,N^2-1$, of such matrices,
the elements $X\in\sun$ are given by
$X=X^aT^a$ with real components $X^a$
(repeated group indices are automatically summed over).
The structure constants $f^{abc}$ in
the commutator relation
\begin{equation}
  [T^a,T^b]=f^{abc}T^c
  \label{eq:A.1}
\end{equation}
are real and totally anti-symmetric in the indices if
the normalization condition
\begin{equation}
  \tr\{T^aT^b\}=-\frac{1}{2}\delta^{ab}
  \label{eq:A.2}
\end{equation}
is imposed.
Moreover, $f^{acd}f^{bcd}=N\delta^{ab}$.

Gauge fields in the continuum theory
take values in the Lie algebra of the gauge group.
Lorentz indices $\mu,\nu,\ldots$ are automatically
summed over when they occur in matching pairs.
The space-time metric is assumed to be euclidean. In particular,
$p^2=p_{\mu}p_{\mu}$ for any momentum $p$ and $\delta_{\mu\mu}=D$.

\section{Flow vertices}\label{app:B}

The flow vertices $\vbulk{2,0},\vbulk{3,0}$ and $\vbulk{1,1}$
are defined through eqs.~\eqref{eq:2.13} and \eqref{eq:6.12}, respectively.
They are explicitly given by
\begin{align}
  \vbulk{2,0}(p,q,r)^{abc}_{\mu\nu\rho}=
  if^{abc}\bigl\{(r-q)_{\mu}\delta_{\nu\rho}
  +2q_{\rho}\delta_{\mu\nu}-2r_{\nu}\delta_{\mu\rho}
  \notag\\[2.5ex]
  +(\alpha_0-1)
  (q_{\nu}\delta_{\mu\rho}-r_{\rho}\delta_{\mu\nu})\bigr\},
  \label{eq:B.1}\\[3.5ex]
  \vbulk{3,0}(p,q,r,s)^{abcd}_{\mu\nu\rho\sigma}=
  f^{abe}f^{cde}
  (\delta_{\mu\sigma}\delta_{\nu\rho}-\delta_{\mu\rho}\delta_{\sigma\nu})
  \notag\\[2.5ex]
  \quad
  +f^{ade}f^{bce}
  (\delta_{\mu\rho}\delta_{\sigma\nu}-\delta_{\mu\nu}\delta_{\rho\sigma})+
  f^{ace}f^{dbe}
  (\delta_{\mu\nu}\delta_{\rho\sigma}-\delta_{\mu\sigma}\delta_{\nu\rho}),
  \label{eq:B.2}\\[3.5ex]
  \vbulk{1,1}(p,q,r)^{abc}_{\mu}=\alpha_0if^{abc}r_{\mu}.
  \label{eq:B.3}
\end{align}
In the Feynman diagrams, the momenta in these formulae are identified
with the ingoing line momenta.

\section{Evaluation of the diagrams in fig.~\ref{fig:brsdiag}}\label{app:C}

The contributions of the diagrams in fig.~\ref{fig:brsdiag} to the
Fourier transform of the correlation functions
in the BRS identity \eqref{eq:6.23} are of the form
\begin{equation}
   -{g_0^4if^{abc}\over \lambda_0 p^2q^2r^2}\mathcal{C}_{\mu,n},
   \label{eq:C.1}
\end{equation}
where $n$ is the diagram number. Explicitly, one obtains
\begin{align}
   \mathcal{C}_{\mu,1}=
   \frac{1}{2}p^2(q-r)_{\mu}\rme^{-t\alpha_0(q^2+r^2)}+
   \frac{1}{2}(q^2-r^2)p_{\mu}\rme^{-t\alpha_0p^2}
   -\mathcal{C}_{\mu,2},
   \label{eq:C.2}\\[3pt]
   \mathcal{C}_{\mu,2}=
   \frac{1}{2}\bigl\{p^2(q-r)_{\mu}+(q^2-r^2)p_{\mu}\bigr\}\rme^{-tp^2},
   \label{eq:C.3}
\end{align}
for the diagrams contributing to the left-hand side of the equation
and
\begin{align}
   \mathcal{C}_{\mu,3}=-p^2r_{\mu}\rme^{-t\alpha_0(q^2+r^2)},
   \label{eq:C.4}\\[3pt]
   \mathcal{C}_{\mu,4}=-\frac{1}{2}p^2p_{\mu}\bigl\{
   \rme^{-t\alpha_0(q^2+r^2)}-\rme^{-t\alpha_0p^2}\bigr\},
   \label{eq:C.5}\\[3pt]
   \mathcal{C}_{\mu,5}=-(qr+r^2)p_{\mu}\rme^{-t\alpha_0p^2},
   \label{eq:C.6}
\end{align}
for the diagrams on the right-hand side. The BRS symmetry requires
\begin{equation}
   \mathcal{C}_{\mu,1}+\mathcal{C}_{\mu,2}=
   \mathcal{C}_{\mu,3}+\mathcal{C}_{\mu,4}+\mathcal{C}_{\mu,5},
   \label{eq:C.7}
\end{equation}
which is indeed the case.

\begin{thebibliography}{7}

% Wilson flow

\bibitem{WilsonFlow}
M. L\"uscher,
{\itshape Properties and uses of the Wilson flow in lattice QCD},
JHEP 1008 (2010) 071

\bibitem{Villasimius}
M. L\"uscher,
{\itshape Topology, the Wilson flow and the HMC algorithm},
XXVIII International Symposium on Lattice Field Theory, June 14-19 2010, Villasimius, Italy,
PoS(Lattice 2010)015


% Renormalization of the Langevin equation

\bibitem{Zinn}
J. Zinn--Justin,
{\itshape Renormalization and stochastic quantization},
Nucl. Phys. B275 [FS17] (1986) 135

\bibitem{ZinnZwanziger}
J. Zinn--Justin, D. Zwanziger,
{\itshape Ward identities for the stochastic quantization of gauge fields},
Nucl. Phys. B295 [FS21] (1988) 297


% BRS symmetry

\bibitem{BRSI}
C. Becchi, A. Rouet and R. Stora,
{\itshape Renormalization of the abelian Higgs-Kibble model},
Comm. Math. Phys. 42 (1975) 127

\bibitem{BRSII}
C. Becchi, A. Rouet and R. Stora,
{\itshape Renormalization of gauge theories},
Ann. Phys. (NY) 98 (1976) 287


% Renormalization of quantum field theories on a half space

\bibitem{Symanzik}
K. Symanzik,
{\itshape Schr\"odinger representation and Casimir effect in
renormalizable quantum field theory},
Nucl. Phys. B190 [FS3] (1981) 1

\end{thebibliography}
